\chapter*{Conclusion générale et Perspectives}
\addcontentsline{toc}{chapter}{Conclusion générale et Perspectives}
\markboth{Conclusion générale et Perspectives}{}

Ce projet a abouti à la conception et au déploiement du BTE Security AI Agent, un
moteur DevSecOps autonome reposant sur une intelligence artificielle locale. Le point
de départ était un VPS Ubuntu vierge ; la plateforme livrée compte onze conteneurs de
production, auxquels s'ajoute une douzième \textit{sandbox} LocalAI, activable à la
demande pour évaluer des moteurs d'inférence concurrents. L'objectif fixé au départ a
été atteint, puis vérifié sur une \textit{Pull Request} de démonstration instrumentée
en fin de stage (PR~\#18) ainsi que sur plusieurs exécutions intermédiaires du pipeline
au fil des sprints. Il consistait à automatiser la détection des vulnérabilités
dans les \textit{Pull Requests} en 15 à 25 minutes environ, sans externaliser le code
source et en couvrant l'intégralité de l'OWASP Top 10.

Exécutés tous deux localement via Ollama, les modèles \texttt{qwen2.\allowbreak{}5-coder:7b}
(classification rapide) et \texttt{qwen2.\allowbreak{}5-coder:14b} (analyse de sécurité
approfondie) ont permis de tenir les exigences de confidentialité bancaire sans sacrifier
la qualité des revues. Fusionner en un seul appel les deux requêtes LLM jusque-là
séquentielles a réduit de moitié la durée du pipeline : d'une fourchette de trente à
cinquante minutes, elle est tombée à quinze ou vingt-cinq minutes. Nettoyer les sorties
SAST avant de les transmettre au modèle, soit 52\% de tokens en moins, a libéré de la
fenêtre de contexte et amélioré les analyses. Le système anti-hallucination à six couches
a, par ailleurs, mis fin à la fabrication de métriques d'infrastructure dans l'assistant
de chat.

Plusieurs briques techniques concourent à la sécurité, à la fiabilité et à
l'observabilité de la plateforme. L'orchestration LangGraph, adossée à une persistance
PostgreSQL, garantit que le pipeline survit aux redémarrages ; le circuit breaker et les
\textit{fallbacks} le maintiennent opérationnel lorsqu'un service devient ponctuellement
indisponible. Côté supervision, quinze règles d'alerte Prometheus, vingt-huit métriques
personnalisées et trois \textit{dashboards} Grafana donnent une vue complète de
l'infrastructure ; quinze de ces métriques ont été ajoutées au sprint 8.1 pour suivre la
consommation par conteneur, la santé de la sandbox LocalAI et la télémétrie A/B du routeur
de chat. Quant à la méthodologie Scrumban, qui combine la cadence de sprints à durée fixe
héritée de Scrum et une discipline Kanban plafonnée à un seul élément en cours (WIP de 1),
elle a permis, sur dix sprints, d'absorber les imprévus de production : l'urgence disque,
l'AlertManager silencieusement cassé depuis le déploiement, la panne de VictoriaMetrics
ou la réécriture du parser de revue PR. Chacun de ces incidents est devenu la carte
prioritaire du sprint en cours et a donné lieu à un correctif déployé en production avant
la clôture de l'itération.

Plusieurs chantiers restent néanmoins ouverts. Le premier concerne le durcissement de la
surface exposée : les ports hôte de Redis, Prometheus, AlertManager, VictoriaMetrics et
de l'agent sont toujours ouverts et doivent être refermés, Redis attend encore son mot de
passe, et le passage à HTTPS via Let's Encrypt demeure un prérequis pour un usage bancaire.
La résilience opérationnelle vient ensuite : PostgreSQL héberge la base de connaissances
sécurité et les checkpoints LangGraph, mais ne bénéficie d'aucune sauvegarde quotidienne,
et sept conteneurs tournent encore sans \textit{healthcheck} Docker. Reste enfin une
question propre au cœur IA de la plateforme : le diff analysé étant une entrée non fiable,
faire reposer le verdict de blocage sur les \textit{findings} déterministes des scanners,
et non sur le seul jugement du modèle, mettrait la revue à l'abri d'une éventuelle
injection d'instructions dans le code soumis. Sur le terrain de l'inférence,
l'évaluation d'un second moteur conduite en post-Sprint 8 a confirmé le choix d'Ollama,
plus rapide de 22\,\% que LocalAI sur un matériel et un modèle identiques. La sandbox
LocalAI conserve toutefois son intérêt : elle ouvre l'accès à des modèles absents de la
bibliothèque Ollama, tel Qwen3-Coder MoE 30B/3B-actifs, et un \textit{fine-tuning} sur
des CVE bancaires reste une piste à explorer.

L'évolution la plus déterminante pour la suite serait la \textbf{migration vers une
inférence GPU}. Sur le matériel actuel (CPU Haswell 12 cœurs, sans GPU), une revue de
\textit{Pull Request} demande 15 à 25 minutes selon sa complexité, ce qui correspond à un
débit d'environ 3 tokens par seconde pour le modèle \texttt{qwen2.\allowbreak{}5-coder:14B}.
Un GPU haut de gamme (NVIDIA RTX 4090 ou A100) porterait ce débit à 30 ou 50 tokens par
seconde et diviserait donc la durée par un facteur de 5 à 10, ramenant chaque revue à 2 ou
5 minutes. Le coût reste contenu : environ 3\,000~\texteuro{} pour une RTX 4090, ou
500~\texteuro/mois pour du cloud GPU à la demande. Il s'inscrit dans l'enveloppe de la DSI
et constitue le moyen le plus direct d'étendre l'agent à l'ensemble des dépôts BTE sans
toucher à l'architecture logicielle. L'inférence resterait sur l'infrastructure de la
banque, ce qui préserve l'exigence de confidentialité de la BCT, tandis que la revue
deviendrait quasi instantanée pour le développeur.

À plus longue échéance, l'agent a vocation à couvrir l'ensemble des dépôts de la BTE, et
les revues accumulées à alimenter un tableau de bord de risques destiné au RSSI.

Sur le plan personnel, ce stage aura été l'occasion de mettre en pratique l'automatisation
de la sécurité d'une chaîne CI/CD, le déploiement et l'exploitation de grands modèles de
langage en local, et la construction d'un système d'observabilité capable de repérer ses
propres pannes. Plusieurs des difficultés rencontrées (l'urgence disque, le parser de revue
PR à réécrire, le bug de cold-load de LocalAI) ne seraient jamais apparues sans une mise en
conditions réelles ; s'il fallait retenir une leçon de ce travail, ce serait sans doute
celle-là. La plateforme remise à la BTE dépasse le cadre du stage : elle met entre les mains
de la banque un outil concret pour appuyer sa transformation DevSecOps.
